\documentclass[12pt]{article}
\usepackage[portuguese]{babel}
\usepackage[utf8]{inputenc}
% \usepackage{fontspec}
\usepackage{graphicx}
\usepackage{fancyhdr}
\usepackage{extramarks}
\usepackage{xparse}
\usepackage{geometry}
\usepackage{indentfirst}
\usepackage{multicol}
\usepackage{multirow}
\usepackage{wrapfig}
\usepackage{stackengine}
\usepackage{float}
\usepackage{titlesec}
\usepackage[table]{xcolor}
\usepackage{tikz}
\usepackage{amsmath}
\usepackage{lipsum}
\usepackage{array}

%----------------------------------- Fonte -------------------------------------
% \setsansfont{Verdana}
% \usepackage{arev}
\usepackage{dejavu}
\usepackage[T1]{fontenc}
\hyphenpenalty 10000
%\exhyphenpenalty 10000
\renewcommand*\familydefault{\sfdefault}

%-------------------------------- Cor da Fonte ---------------------------------
\definecolor{mediumgray}{rgb}{0.35, 0.35, 0.35}
\definecolor{BlueEd}{rgb}{0,0,102}
\definecolor{OrangeEd}{rgb}{1.0,0.45,0.30}

\makeatletter
\newcommand{\globalcolor}[1]{%
  \color{#1}\global\let\default@color\current@color
}
\makeatother

\AtBeginDocument{\globalcolor{mediumgray}}

%------------------------------- Formatação ------------------------------------
\let\CLASSsection\section

\RenewDocumentCommand{\section}{som}{%
  \IfBooleanTF{#1}
   {% there's a *
    \CLASSsection*{#3}\markboth{#3}{}%
   }
   {% no *
    \IfNoValueTF{#2}
     {% no opt arg
      \CLASSsection{#3}%
     }
     {% opt arg
      \CLASSsection[#2]{#3}%
     }%
   }%
}

% \geometry{
%     a4paper,
%     total={160mm,257mm},
%     left=25mm,
%     top=30mm,
%     headheight=15mm,
%     includeheadfoot,
%     voffset=-1in,
%     hoffset=0in
% }

\geometry{
    a4paper,
    %total={160mm,257mm},
    left=20mm,
    right=20mm,
    top=6mm,
    bottom=19mm,
    headheight=16mm,
    includeheadfoot,
    %voffset=-1in,
    hoffset=0in
}
 
%---------------------------------- Cabeçalho ----------------------------------

\newcolumntype{?}{!{\vrule width 1.5pt}}
%\renewcommand{\arraystretch}{1}

\newcommand{\sectionauthor}{Autor da Seção}
\newcommand{\sectionclass}{Disciplina}
\renewcommand{\headrulewidth}{0pt}
\pagestyle{fancy}
\renewcommand{\sectionmark}[1]{\markboth{\thesection\ #1}{}}
\fancyhf{}

\fancyhead[L]{ %------ Campo da esquerda ------
    \begin{tabular}{@{}c@{}}
        \includegraphics[width=20mm]{assets/logo.png}
    \end{tabular}
    \color{OrangeEd}
    \fontsize{8pt}{11}\selectfont
    %\renewcommand{\arraystretch}{0.7}
    \begin{tabular}{l@{}l@{}}
        \mbox{}\\
        \multicolumn{1}{l?}{} & \quad \textbf{E-book didático}\\ 
        \multicolumn{1}{l?}{} & \quad \textbf{Enem e Vestibular}\\
        \mbox{}
    \end{tabular}
}

\fancyhead[R]{  %------ Campo da direita ------
    %\renewcommand{\arraystretch}{1}
    \fontsize{10pt}{12}\selectfont
    \begin{tabular}{@{}r@{}}
        \color{OrangeEd}\MakeUppercase{\textbf{\sectionclass}} \\
        \textbf{\firstleftmark} \\
        %Prof. Matheus Silva
        \sectionauthor
    \end{tabular}
}


%----------------------------------- Rodapé ------------------------------------
\fancyfoot[L]{  %------ Campo da esquerda ------
    \color{black}
    \fontsize{8pt}{8}\selectfont Educursinho: Educação e aprovação  \quad\textbf{|}\quad 2021
}

\fancyfoot[R]{  %------ Campo da direita ------
    \color{black}
    \fontsize{8pt}{8}\selectfont \thepage
}
%--------------------------- Fim Rodapé e Cabeçalho ----------------------------


%-------------------------------- Formatação -----------------------------------
\titleformat*{\section}{\color{OrangeEd}\fontsize{18pt}{18}\centering\bfseries\sffamily\uppercase}
\titleformat*{\subsection}{\color{black}\raggedright\large\bfseries\sffamily\uppercase}
\titleformat*{\subsubsection}{\color{black}\bfseries\uppercase}
% \titlespacing{\subsection}{0mm}{0mm}{0mm}
% \titlespacing{\section}{0mm}{0mm}{0mm}

\titlespacing{\subsection}{0mm}{5mm}{8mm}
\titlespacing{\section}{0mm}{0mm}{8.5mm}

\setlength{\columnsep}{12.5mm}
\setlength{\parindent}{12.5mm} 
%\setlength{\parskip}{1mm}
\renewcommand{\baselinestretch}{1.15}

%------------------------------ Inicio documento -------------------------------    
\begin{document}
    
    \renewcommand{\sectionclass}{Matemática}

    %\renewcommand{\sectionauthor}{Prof. Rafael Simões}

\begin{multicols*}{2}
        
    %\section*{Gráficos e Funções}
    \section*{Eletroquímica I}
    \lipsum[1]

    \subsection*{Funções algébricas}
    \lipsum[1-2]
    
    \noindent \\Resolução:
    \begin{align}
        f(x) &= \frac{sen(x)}{cos(x)} \\[10pt]
        f(x) &= tan(x)
    \end{align}
    

    %\pagebreak
    \newpage
    \section*{Funções polinomiais}
    \lipsum[1]

    \subsection*{Raízes}
    \lipsum[1-2]    
    
\end{multicols*}
    \renewcommand{\sectionauthor}{Gustavo Ale}

\begin{multicols*}{2}

    \section*{Geometria}
    A geometria é provavelmente a área mais antiga da matemática, precursora da própria álgebra e sendo um dos pilares
    da matemática. A palavra geometria é resultado da combinação das palavras gregas \textit{geo} e
    \textit{metron}/\textit{metri}, significando respectivamente Terra e medição, pois a área da geometria consiste da
    medição e entendimento das relações e propriedades  contidas nas fíguras geométricas: comprimento, distância, ângulo,
    área, volume e perímetro, as fíguras geométricas por sua vez são parte integrante da natureza e da Terra (\textit{geo}).

    \subsection*{Definições gerais}
    Antes do estudo da geometria devemos ter ciência das definições de conceitos presentes nesse ramo da matemática,
    dentro destes conceitos se enquadram as medidas, os ângulos, as fíguras geométricas e suas componentes.
    
    \vfill\null
    \columnbreak
    \vfill\null

    \subsection*{Grandezas e unidades de medida}
    \subsubsection*{Distância e comprimento}
    Distância e comprimento medem o quão longe dois pontos estão entre si, e no caso de arestas essa medida é
    denotada como comprimento. A unidade de medida base utilizada para comprimento é o metro (símbolo ${m}$), além do
    metro existem seus múltiplos que são igualmente utilizados dependendo da distância/comprimento aferido.
    \begin{table}[H]
        \resizebox{\columnwidth}{!}{%
            \begin{tabular}{l|l|l}
                \rowcolor{gray!25}
                \textbf{Nome} & \textbf{Sigla} & \textbf{Equivalência} \\ \hline
                picometro     & $pm$           & $10^{-12}m$           \\ \rowcolor{gray!10}
                nanometro     & $nm$           & $10^{-9}m$            \\
                micrometro    & $\mu m$        & $10^{-6}m$            \\ \rowcolor{gray!10}
                milimetro     & $mm$           & $10^{-3}m$            \\
                centímetro    & $cm$           & $10^{-2}m$            \\ \rowcolor{gray!10}
                decímetro     & $dm$           & $10^{-1}m$            \\
                metro         & $m$            & $1m$                  \\ \rowcolor{gray!10}
                decâmetro     & $dam$          & $10m$                 \\
                hêctometro    & $hm$           & $100m$                \\ \rowcolor{gray!10}
                quilometro    & $km$           & $1000m$
            \end{tabular}%
        }
        %\caption{Unidades medida de comprimento e distância}
    \end{table}
    Outras medidas de comprimento são: \\
    \textbf{Perímetro:} comprimento do contorno de uma fígura geométrica. \\
    \textbf{Raio:} distância entre o centro de uma circunferência até seu contorno ou superfície.\\
    \textbf{Diâmetro:} comprimento de reta que passe pelo centro da circunferência e cujo seus pontos de
    início e fim estejam sobre a circunferência\\
    \textbf{Circunferência:} perímetro de uma circunferência.\\

    \begin{figure}[H]
        \centering
        %\resizebox{\columnwidth}{!}{%
            \begin{tikzpicture}
                \coordinate (A) at (2,2);
                \coordinate (B) at (4,2);
                \coordinate (C) at (2,0);
                \coordinate (D) at (2,4);
                \draw (A)--node[below] {$r$}(B) ;
                \draw (C)--node[left] {$d$}(D) ;
                \draw (A) circle (2cm)
                (0,3) node[left] {$c$};
            \end{tikzpicture}
        %}
        \caption{Raio $r$, diâmetro $d$ e circunferência $c$ de um círculo}
    \end{figure}

    % \begin{figure}[H]
    %     \centering
    %     \includegraphics[width=0.7\columnwidth]{assets/image_placeholder.png}
    %     \caption{Raio, diâmetro e circunferência de um círculo}
    % \end{figure}

    \subsubsection*{Área}
    Área é a medida que expressa a quantidade de espaço bidimenssional ocupado por uma fígura geométrica,
    área de superfície é o equivalente da área para uma superfície ou face de um objeto tridimenssional. A unidade base
    para área é o metro quadrado (símbolo ${m^2}$). Seus múltiplos também acompanham o termo 'quadrado',
    e a equivalência é a mesma do metro, porém elevada ao quadrado. Ex.: $1m = 10^2cm$ e $1m^2 = (10^2cm)^2 = 10^4cm^2$,
    lembrar essa regra pode facilitar na hora da conversão de área.

    \begin{table}[H]
        \resizebox{\columnwidth}{!}{%
            \begin{tabular}{l|l|l}
                \rowcolor{gray!25}
                \textbf{Nome}       & \textbf{Sigla} & \textbf{Equivalência}       \\ \hline
                picometro quadrado  & $pm^2$         & $10^{-24}m^2$               \\ \rowcolor{gray!10}
                nanometro quadrado  & $nm^2$         & $10^{-18}m^2$               \\
                micrometro quadrado & $\mu m$        & $10^{-12}m^2$               \\ \rowcolor{gray!10}
                milimetro quadrado  & $mm^2$         & $10^{-6}m^2$                \\
                centímetro quadrado & $cm^2$         & $10^{-4}m^2$                \\ \rowcolor{gray!10}
                decímetro quadrado  & $dm^2$         & $10^{-2}m^2$                \\
                metro quadrado      & $m^2$          & $1m^2$                      \\ \rowcolor{gray!10}
                decâmetro quadrado  & $dam^2$        &                             \\ \rowcolor{gray!10}
                are                 & $are$          & \multirow{-2}{*}{$100m^2$}  \\
                hêctometro quadrado & $hm^2$         &                             \\
                hectare             & $ha$           & \multirow{-2}{*}{$10^4m^2$} \\ \rowcolor{gray!10}
                quilometro quadrado & $km^2$         & $10^6m^2$
            \end{tabular}%
        }
        %\caption{Unidades medida de área}
    \end{table}

    \subsubsection*{Volume}
    Assim como a área expressa a quantidade de espaço bidimenssional ocupado por uma fígura geométrica, o volume expressa
    a quantidade de espaço tridimenssional ocupado por um objeto. No caso do volume existe duas unidades de medida base,
    o metro cúbico e o litro, ambos são usados regularmente, na qual o litro é a unidade mais usada para volumes
    pequenos, enquanto o metro cúbico é usado para expressar volumes maiores.

    \begin{table}[H]
        \resizebox{\columnwidth}{!}{%
            \begin{tabular}{l|l|l}
                \rowcolor{gray!25}
                \textbf{Nome}     & \textbf{Sigla} & \textbf{Equivalência}          \\ \hline
                milimetro cúbico  & $mm^3$         & $10^{-9}m^3$                   \\ \rowcolor{gray!10}
                centímetro cúbico & $cm^3$         &                                \\ \rowcolor{gray!10}
                mililitro         & $ml$           & \multirow{-2}{*}{$10^{-6}m^3$}  \\
                decímetro cúbico  & $dm^3$         &                                \\ %\hline
                litro             & $l$            & \multirow{-2}{*}{$10^{-3}m^3$} \\ \rowcolor{gray!10}
                metro cúbico      & $m^3$          & $1m^2$                         \\ %\hline                      
            \end{tabular}%
        }
        %\caption{Unidades medida de volume}
    \end{table}

    \subsubsection*{Ângulo}

    Ângulo é uma medida de inclinação entre duas retas ou dois planos, desde que eles não sejam paralelos entre si, no 
    contexto da geometria os ângulos estão presentes em todos os vértices das fíguras geométricas. 
    As unidades de medida mais comuns para ângulo são graus (símbolo $^{\circ}$), radianos (símbolo $rad$) e gradianos (símbolo $gon$).
    Quanto no escopo das funções trigonométricas a unidade de medida mais utilizada é o radiano, por outro lado, 
    em aplicações de engenharia e na geometria a unidade mais usada é o grau. Considerando um grau $\alpha$ em graus, 
    sua conversão para radianos se dá por $rad(\alpha) = \alpha\cdot\pi/180$ 

    % teste matemática em duas colunas
    % \color{black}
    % \begin{align}
    %     f(x) &= \frac{sen(x)}{cos(x)} \\[10pt]
    %     f(x) &= tan(x)
    % \end{align}
    % \color{mediumgray}

    %\pagebreak
    % \newpage
    % \section*{Ângulos}
    % \lipsum[1]

    % \subsection*{Simetrias de figuras planas}
    % \lipsum[1-2]    

    \subsection*{Fíguras geométricas e suas componentes}

    \subsubsection*{Ponto}
    A primeira componente das fíguras geométricas e a mais simples delas é o ponto, o mesmo deve ser imaginado como um
    ponto infinitesimal, não possuindo comprimento, área ou perímetro. Quando um ponto é integrante de uma fígura
    geométrica o mais comum é chamarmos ele de \textit{vértice}
    \begin{figure}[H]
        \centering
        %\resizebox{\columnwidth}{!}{%
            \begin{tikzpicture}
                \coordinate (A) at (1,1);
                \coordinate (B) at (5,3);
                \coordinate (C) at (5,0);
                \node[circle, fill, label={above:$A$}, inner sep=2pt] at (A) {};
                \node[circle, fill, label=, inner sep=2pt] at (B) {};
                \node[circle, fill, label=, inner sep=2pt] at (C) {};
                \draw (A)--(B);
                \draw (B)--(C);
                \draw (A)--(C);
            \end{tikzpicture}
        %}
        \caption{Triângulo com o vértice $A$ em realce }
        % \label{fig:vertice_tri}
    \end{figure}


    % \begin{figure}[H]
    %     \centering
    %     \includegraphics[width=0.7\columnwidth]{assets/image_placeholder.png}
    %     \caption{Vértice de um triângulo }
    % \end{figure}

    A partir de dois pontos $A$ e $B$ podemos traçar uma reta $\overline{AB}$ entre eles e caso exista um terceiro ponto $C$ sob a mesma reta, 
    então esses três pontos são \textbf{colineares}.    

    \begin{figure}[H]
        \centering
        %\resizebox{\columnwidth}{!}{%
            \begin{tikzpicture}
                \coordinate (A) at (0,0);
                \coordinate (B) at (5,2);
                \coordinate (C) at ($(A)!0.5!(B)$);
                \node[circle, fill, label={above:$A$}, inner sep=2pt] at (A) {};
                \node[circle, fill, label={above:$B$}, inner sep=2pt] at (B) {};
                \node[circle, fill, label={above:$C$}, inner sep=2pt] at (C) {};
                \draw (A)--(B);
            \end{tikzpicture}
        %}
        \caption{Pontos colineares $A$,$B$ e $C$}
        \label{fig:reta_ab}
    \end{figure}

    % \begin{figure}[H]
    %     \centering
    %     \includegraphics[width=0.7\columnwidth]{assets/image_placeholder.png}
    %     \caption{Pontos colineares}
    %     \label{fig:reta_abc}
    % \end{figure}

    \subsubsection*{Reta}
    A reta é um conjunto de pelo menos dois pontos, no caso dois pontos quaisquer $A$ e $B$ podem gerar a reta $\overline{AB}$, como na fígura \ref{fig:reta_ab}. 
    Assim como o ponto é idealmente infinitesimal a reta idealmente não possui espessura. Quando uma reta é parte 
    integrante de uma fígura geométrica sua denominação é \textit{aresta}.

    \begin{figure}[H]
        \centering
        %\resizebox{\columnwidth}{!}{%
            \begin{tikzpicture}
                \coordinate (A) at (0,0);
                \coordinate (B) at (3,3);
                \coordinate (C) at (4,1);
                \node[circle, fill, label={above:$A$}, inner sep=2pt] at (A) {};
                \node[circle, fill, label={above:$B$}, inner sep=2pt] at (B) {};
                \node[circle, fill, label={right:$C$}, inner sep=2pt] at (C) {};
                \draw (A)--(B);
                \draw (B)--(C);
                \draw (A)--(C);
            \end{tikzpicture}
        %}
        \caption{Retas $\overline{AB}$,$\overline{BC}$ e $\overline{AC}$ do triângulo $ABC$}
        \label{fig:tri_abc}
    \end{figure}

    % \begin{figure}[H]
    %     \centering
    %     \includegraphics[width=0.7\columnwidth]{assets/image_placeholder.png}
    %     \caption{Reta $\overline{AB}$ do triângulo $\overline{ABC}$}
    %     \label{fig:reta_ab}
    % \end{figure}

    Uma reta $p$ é \textbf{perpendicular} a reta $q$ se o ângulo formado entre elas for de 90 graus.
    
    \begin{figure}[H]
        \centering
        %\resizebox{\columnwidth}{!}{%
            \begin{tikzpicture}
                \coordinate (O) at (1,1.5);
                \coordinate (A) at (0,2);
                \coordinate (B) at (4,0);
                \coordinate (C) at (0,0);
                \coordinate (D) at (3,4.5);
                \draw (A)--(B) node[above] {$p$}  ;
                \draw (C)--(D) node[left] {$q$};
                % \pic[angle radius=1cm,"$\alpha$"] {angle=a--b--c};
                \pic ["$90^{\circ}$", draw, -, angle radius=1cm] {angle = B--O--D};
                % \draw (B)--(C);
                % \draw (A)--(C);
            \end{tikzpicture}
        %}
        \caption{Retas perpendiculares $p$ e $q$}
    \end{figure}

    % \begin{figure}[H]
    %     \centering
    %     \includegraphics[width=0.7\columnwidth]{assets/image_placeholder.png}
    %     \caption{Retas perpendiculares $p$ e $q$}
    % \end{figure}

    %\columnbreak
    Duas retas $p$ e $q$ são consideradas \textbf{paralelas entre si}, se qualquer reta $r$ perpendicular a reta $p$ também for 
    perpendicular a reta $q$. Outra forma de descrever retas paralelas é através do \textbf{Quinto Postulado de Euclides}
    que diz: Supondo que duas retas $p$ e $q$ são cortadas por uma terceira reta $r$. Se a soma dos ângulos formados um 
    mesmo lado da reta $r$ resultar em 180 graus, então $m$ e $n$ são retas paralelas entre si.

    \begin{figure}[H]
        \centering
        %\resizebox{\columnwidth}{!}{%
            \begin{tikzpicture}
                \coordinate (A) at (0,0);
                \coordinate (B) at (5,1);
                \coordinate (C) at (0,2);
                \coordinate (D) at (5,3);
                \draw (A)--node[below] {$p$}(B) ;
                \draw (C)--node[above] {$q$}(D) ;
            \end{tikzpicture}
        %}
        \caption{Retas paralelas $p$ e $q$}
    \end{figure}

    % \begin{figure}[H]
    %     \centering
    %     \includegraphics[width=0.7\columnwidth]{assets/image_placeholder.png}
    %     \caption{Retas paralelas $p$ e $q$}
    % \end{figure}


    \subsubsection*{Plano e superfície}
    Quando se tem 3 ou mais pontos o resultado da conexão deles é tido como superfície, plano ou fígura geométrica
    dependendo da área da matemática e do contexto de estudo, porém vale ter em mente que também existem superfícies
    não bidimensionais. Uma excessão a isso são os círculos e as elipses que não possuem vértices. No contexto da
    geometria a ser estudado nesse capítulo a conexão de 3 ou mais vértices geram exclusivamente fíguras geométricas.

    \begin{figure}[H]
        \centering
        %\resizebox{\columnwidth}{!}{%
            \begin{tikzpicture}
                \coordinate (A) at (0,0);
                \coordinate (B) at (4,1);
                \coordinate (C) at (5,4);
                \coordinate (D) at (2,3);
                \draw (A)-- node[below] {$\overline{AB}$} (B) ;
                \draw (B)--(C) ;
                \draw (C)-- node[above] {$\overline{CD}$} (D) ;
                \draw (D)--(A) ;
                \node[circle, fill, label={left:$A$}, inner sep=2pt] at (A) {};
                \node[circle, fill, label={right:$B$}, inner sep=2pt] at (B) {};
                \node[circle, fill, label={right:$C$}, inner sep=2pt] at (C) {};
                \node[circle, fill, label={left:$D$}, inner sep=2pt] at (D) {};
            \end{tikzpicture}
        %}
        \caption{Um quadrilátero $ABCD$ com as arestas $\overline{AB}$ e $\overline{CD}$ realçadas}
    \end{figure}

    % \begin{figure}[H]
    %     \centering
    %     \includegraphics[width=0.7\columnwidth]{assets/image_placeholder.png}
    %     \caption{Um quadrilátero com as arestas realçadas em vermelho}
    % \end{figure}

    \subsubsection*{Sólidos}
    No contexto de fíguras geométricas tridimenssionais, ou chamados sólidos, os sólidos podem ser formados por 4 ou
    mais pontos ou através manipulação de fíguras geométricas bidimensionais num espaço tridimenssional.
    Todos objetos físicos que conhecemos podem ser abstraídos como formas geométricas, independente da complexidade.

    %\columnbreak
    Grandezas notórias das componentes citadas:
    \begin{itemize}
        \setlength\itemsep{1.15pt}
        \item Ponto, vértice: distância relativo a outro objeto.
        \item Reta: comprimento, distância e ângulo relativos a outro objeto.
        \item Superfície: área, perímetro.
        \item Sólido: volume.
    \end{itemize}

    \begin{figure}[H]
        \centering
        \includegraphics[width=0.7\columnwidth]{assets/image_placeholder.png}
        \caption{Tetraedro, o objeto tridimenssional com menos faces}
    \end{figure}

    

\end{multicols*}
    %\include{Parte3}
    
\end{document}